% TOP_PROBLEM: {"problemId": "problem.grothendiecks-section-conjecture-in-anabelian-geometry", "canonicalTitle": "Grothendiecks Section Conjecture in Anabelian Geometry", "releaseRank": 100, "primaryDomain": "number_theory_arithmetic_geometry", "primaryDomainLabel": "Number theory & arithmetic geometry", "displayStatus": "open", "releaseStatus": "open"}
% !TeX program = lualatex
% Definition and sources only; this file is not an LLM solution attempt.
\documentclass[11pt]{article}
\usepackage[margin=25mm]{geometry}
\usepackage{fontspec}
\setmainfont{Latin Modern Roman}
\usepackage{amsmath,amssymb,mathtools}
\usepackage{unicode-math}
\setmathfont{Latin Modern Math}
\usepackage{xurl,hyperref}
\hypersetup{colorlinks=true,urlcolor=blue}
\setcounter{secnumdepth}{0}
\setlength{\parindent}{0pt}
\setlength{\parskip}{5pt}
\setlength{\emergencystretch}{3em}
\begin{document}
\section*{\texorpdfstring{Grothendiecks Section Conjecture in Anabelian Geometry}{Grothendiecks Section Conjecture in Anabelian Geometry}}\label{100}
\subsection{Definitions and mathematical statement}
For a smooth projective geometrically connected curve $C/K$ of genus at least two, choose a geometric base point and consider
\[
 1\longrightarrow\pi_1(C_{\bar K})\longrightarrow\pi_1(C)
 \longrightarrow G_K\longrightarrow1.
\]
A section is a continuous homomorphism $s:G_K\to\pi_1(C)$ whose composition with the projection is the identity. Sections are identified under conjugation by $\pi_1(C_{\bar K})$. A rational point supplies such a conjugacy class through functoriality of the étale fundamental group. The conjecture says that the map from $C(K)$ to these section classes is bijective.
\par\smallskip\textit{Formulation and concept references:} \href{https://doi.org/10.1017/CBO9780511758874.005}{[S1]}.

\subsection{Short English statement}
Do rational points on a hyperbolic number-field curve correspond exactly to the sections of its arithmetic fundamental-group sequence?
\subsection{Sources}
\begin{itemize}
\item[S1]A. Grothendieck, Brief an G. Faltings, LMS Lecture Notes 242:49-58, 1997.
\url{https://doi.org/10.1017/CBO9780511758874.005}.
\textit{Formulation source.}
\item[S2]Polylogarithmic Chabauty--Kim loci over number fields.
\url{https://arxiv.org/abs/2608.20615}.
\textit{Further reference; not a proof endorsement.}
\item[S3]On the section conjecture over fields of finite type.
\url{https://onlinelibrary.wiley.com/doi/full/10.1002/mana.70049}.
\textit{Further reference; not a proof endorsement.}
\end{itemize}
\end{document}
